\section{Mid-Term Examination --- Set B}
\label{sec:midterm_set_b}

\LPUCSHeader{Mid-Term Examination (Comprehensive MCQ)}{B}{90 Minutes}{30 Marks}{CO1, CO2 \& CO3 (Unit 1: System, Unit 2: OS, Unit 3: Linux)}

\begin{center}
\sffamily\textbf{\color{AlertRed}Instructions: All 30 questions are compulsory. Each question carries 1 mark. A penalty of 0.25 marks is deducted for each incorrect response.}
\end{center}

\vspace{3mm}

% ------------------------------------------------------------------------------
% UNIT 1: COMPUTER SYSTEM (Q1 TO Q10)
% ------------------------------------------------------------------------------
\mcqitem{1}{As one traverses downward through the computer memory hierarchy from CPU registers to secondary storage disks:}{Access speed decreases, cost per bit decreases, and total storage capacity increases}{Access speed increases, cost per bit increases, and capacity decreases}{Latency decreases to zero}{Both cost and speed remain constant}{1}

\mcqitem{2}{The performance metric of a cache memory defined as the fraction of memory accesses satisfied directly by the cache is the:}{Cache Hit Ratio ($H$)}{Bandwidth Ratio}{Clock Multiplier}{Eviction Index}{1}

\mcqitem{3}{In a Write-Through cache configuration, when the processor writes data to a memory location:}{The update is written simultaneously to both the cache and main memory}{The update is held in cache until the block is evicted}{Data is written directly to the SSD without updating cache}{Writes are permanently buffered in registers}{1}

\mcqitem{4}{The physical lifespan and write endurance of a Solid-State Drive (SSD) is standardized and rated in:}{TBW (Terabytes Written) or P/E (Program/Erase) cycles}{Rotational revolutions per minute (RPM)}{Number of optical laser passes}{Total gigahertz cycles}{1}

\mcqitem{5}{In a RAID 1 storage setup comprising two $4\text{ TB}$ hard drives, the level of storage overhead dedicated strictly to redundancy is:}{$50\%$ (data is duplicated 1:1 on the mirror drive)}{$0\%$}{$75\%$}{$25\%$}{1}

\mcqitem{6}{Which RAID configuration uses dual distributed parity ($P+Q$) based on Reed-Solomon error correction to tolerate the concurrent failure of any two disks?}{RAID 6}{RAID 5}{RAID 0}{RAID 1}{1}

\mcqitem{7}{Graphics Processing Units (GPUs) process computations based on which architectural execution model?}{SIMD / SIMT (Single Instruction, Multiple Data / Threads)}{MIMD with out-of-order execution only}{Strict sequential SISD}{Analog pipeline routing}{1}

\mcqitem{8}{In modern PC architectures, the CPU core frequency is generated by multiplying the motherboard base clock ($100\text{ MHz}$) using a hardware:}{Phase-Locked Loop (PLL) clock multiplier}{Capacitive divider}{DRAM refresh timer}{DMA controller}{1}

\mcqitem{9}{Which processor register holds the data word fetched from memory before it is loaded into other CPU registers or sent to the ALU?}{Memory Data Register (MDR)}{Program Counter (PC)}{Stack Pointer (SP)}{Status Register}{1}

\mcqitem{10}{Under Port-Mapped I/O (Isolated I/O), input/output peripheral devices:}{Use a completely separate address space and dedicated assembly instructions (such as \texttt{IN} and \texttt{OUT})}{Share the same address space and memory instructions as physical RAM}{Cannot send hardware interrupts to the CPU}{Bypass system buses entirely}{1}

% ------------------------------------------------------------------------------
% UNIT 2: OPERATING SYSTEM (Q11 TO Q20)
% ------------------------------------------------------------------------------
\mcqitem{11}{Why does the operating system prevent user-mode applications from executing privileged instructions such as \texttt{HLT} (halt) or \texttt{CLI} (disable interrupts)?}{To maintain system stability and prevent rogue user applications from seizing control or crashing the machine}{Because user programs do not have access to an ALU}{To conserve motherboard battery power}{Because privileged instructions only run on graphics cards}{1}

\mcqitem{12}{In a microkernel operating system, if a device driver crashes during runtime:}{Only that isolated user-space driver process fails and can be restarted without crashing the entire operating system}{The entire operating system kernel halts immediately}{All physical RAM contents are permanently corrupted}{The motherboard firmware must be reflashed}{1}

\mcqitem{13}{In process state transitions, a process moves from the 'Ready' state to the 'Running' state when:}{The Short-Term CPU Scheduler selects the process and dispatches it to the CPU}{Its allocated CPU time slice expires}{It executes a blocking file read}{It terminates execution}{1}

\mcqitem{14}{Which preemptive CPU scheduling algorithm always selects the process that has the smallest amount of CPU execution time remaining until completion?}{Shortest Remaining Time First (SRTF)}{First-Come, First-Served (FCFS)}{Round Robin (RR)}{Priority Scheduling without preemption}{1}

\mcqitem{15}{Which of the following is NOT one of the four essential Coffman conditions necessary for a system deadlock?}{Resource starvation with preemptive eviction}{Mutual Exclusion}{Hold and Wait}{Circular Wait}{1}

\mcqitem{16}{Internal fragmentation occurs in memory management when:}{A fixed-size memory block allocated to a process is larger than requested, leaving wasted memory inside the block}{Free memory is fragmented into many tiny unusable gaps outside allocated partitions}{Page tables exceed the size of the hard drive}{Processes share memory using symbolic links}{1}

\mcqitem{17}{When a Translation Lookaside Buffer (TLB) miss occurs during virtual memory address resolution, the CPU must:}{Perform a page table walk in main physical RAM to locate the required frame mapping}{Immediately terminate the faulting process}{Crash the operating system kernel}{Defragment the hard disk drive}{1}

\mcqitem{18}{A statically linked binary executable differs from a dynamically linked binary because the statically linked binary:}{Embeds all required library code directly into the executable file, eliminating external runtime library dependencies}{Requires shared \texttt{.so} files in \texttt{/usr/lib} at execution time}{Consumes less physical disk storage}{Cannot be executed on multi-core processors}{1}

\mcqitem{19}{In the call stack frame (activation record) created during a C function invocation, which critical datum is saved to ensure execution resumes correctly after function return?}{The Return Address (saved Program Counter)}{The entire hard disk partition table}{The global BSS segment}{The MAC address of the network interface}{1}

\mcqitem{20}{Which stage of the software compilation pipeline parses tokens into a hierarchical Abstract Syntax Tree (AST) while verifying language grammar rules?}{Syntax Analysis (Parser)}{Lexical Analysis (Scanner)}{Code Generation}{Linker}{1}

% ------------------------------------------------------------------------------
% UNIT 3: LINUX OPERATING SYSTEM (Q21 TO Q30)
% ------------------------------------------------------------------------------
\mcqitem{21}{The free software licensing model that governs the Linux kernel and ensures that modified derivative versions remain open source is the:}{GNU General Public License (GPL)}{Proprietary Commercial License}{BSD License without attribution}{Creative Commons Zero (CC0)}{1}

\mcqitem{22}{The official package management tool used by Debian, Ubuntu, and their derivatives to download, install, and update \texttt{.deb} packages is:}{\texttt{apt} (Advanced Package Tool)}{\texttt{dnf}}{\texttt{yum}}{\texttt{pacman}}{1}

\mcqitem{23}{According to the Linux Filesystem Hierarchy Standard (FHS), the directory \texttt{/sbin} is designated specifically for:}{Essential system administration command binaries (such as \texttt{fdisk}, \texttt{iptables}, and \texttt{reboot})}{Standard user text files and scripts}{Removable media mount points}{Source code repositories}{1}

\mcqitem{24}{In Linux file permissions, executing \texttt{chmod 644 report.txt} establishes which permission flags for Owner, Group, and Others?}{Owner: \texttt{rw-}, Group: \texttt{r--}, Others: \texttt{r--}}{Owner: \texttt{rwx}, Group: \texttt{r-x}, Others: \texttt{r-x}}{Owner: \texttt{rw-}, Group: \texttt{rw-}, Others: \texttt{rw-}}{Owner: \texttt{r--}, Group: \texttt{---}, Others: \texttt{---}}{1}

\mcqitem{25}{Which Linux command changes the user and group ownership of a specified file or directory?}{\texttt{chown}}{\texttt{chmod}}{\texttt{umask}}{\texttt{touch}}{1}

\mcqitem{26}{If a Linux user environment has a default \texttt{umask} value of \texttt{022}, what will be the default permissions assigned to a newly created regular text file?}{\texttt{644} (\texttt{-rw-r--r--})}{\texttt{755} (\texttt{-rwxr-xr-x})}{\texttt{777} (\texttt{-rwxrwxrwx})}{\texttt{600} (\texttt{-rw-------})}{1}

\mcqitem{27}{The special permission bit SUID (Set User ID, octal \texttt{4000}), when enabled on an executable file like \texttt{/usr/bin/passwd}, enables:}{Any user executing the binary to temporarily run it with the permissions of the file owner (root)}{The file to be permanently write-protected}{The file to be hidden from standard directory listings}{The file to be executed only during system boot}{1}

\mcqitem{28}{In the Bash shell, the pipe operator (\texttt{|}) functions by:}{Directing the standard output (\texttt{stdout}) of the left command into the standard input (\texttt{stdin}) of the right command}{Writing command output directly into the boot sector}{Terminating both commands immediately}{Executing commands in reverse alphabetical order}{1}

\mcqitem{29}{Which Bash command displays a complete, detailed snapshot of all active processes currently executing on the system across all users?}{\texttt{ps aux}}{\texttt{ls -la}}{\texttt{cat /etc/hosts}}{\texttt{df -h}}{1}

\mcqitem{30}{Which category of hypervisor operates as a software application on top of an existing host operating system (e.g., Oracle VM VirtualBox)?}{Type 2 Hypervisor (Hosted)}{Type 1 Hypervisor (Bare-Metal)}{Native Microkernel}{Container Runtime}{1}

\vspace{5mm}
\hrule
\vspace{5mm}

% ------------------------------------------------------------------------------
% ANSWER KEY & EXPLANATIONS
% ------------------------------------------------------------------------------
\subsection*{Answer Key (Mid-Term Set B)}

\begin{center}
\begin{tabular}{|c|c||c|c||c|c||c|c||c|c||c|c|}
\hline
\textbf{Q} & \textbf{Ans} & \textbf{Q} & \textbf{Ans} & \textbf{Q} & \textbf{Ans} & \textbf{Q} & \textbf{Ans} & \textbf{Q} & \textbf{Ans} & \textbf{Q} & \textbf{Ans} \\
\hline
1 & \textbf{A} & 6 & \textbf{A} & 11 & \textbf{A} & 16 & \textbf{A} & 21 & \textbf{A} & 26 & \textbf{A} \\
2 & \textbf{A} & 7 & \textbf{A} & 12 & \textbf{A} & 17 & \textbf{A} & 22 & \textbf{A} & 27 & \textbf{A} \\
3 & \textbf{A} & 8 & \textbf{A} & 13 & \textbf{A} & 18 & \textbf{A} & 23 & \textbf{A} & 28 & \textbf{A} \\
4 & \textbf{A} & 9 & \textbf{A} & 14 & \textbf{A} & 19 & \textbf{A} & 24 & \textbf{A} & 29 & \textbf{A} \\
5 & \textbf{A} & 10 & \textbf{A} & 15 & \textbf{A} & 20 & \textbf{A} & 25 & \textbf{A} & 30 & \textbf{A} \\
\hline
\end{tabular}
\end{center}

\vspace{3mm}
\subsection*{Detailed Technical Solutions}

\begin{solutionbox}[Key Architectural, Systems \& Linux Explanations]
\textbf{Q.5 (Ans: A)} RAID 1 creates an exact bit-for-bit mirror of data across drives. For two drives of size $S$, total physical capacity is $2S$ while usable storage is $1S$. Exactly $50\%$ of physical space is allocated to mirroring redundancy.\\[2mm]
\textbf{Q.14 (Ans: A)} SRTF is the preemptive variant of Shortest Job First. Whenever a new process arrives in the ready queue whose remaining CPU burst is shorter than the remaining execution time of the currently running process, the CPU is preempted and reassigned to the new arrival.\\[2mm]
\textbf{Q.24 (Ans: A)} Octal 6 is $4+2+0 \implies \texttt{rw-}$ (Read + Write). Octal 4 is $4+0+0 \implies \texttt{r--}$ (Read only). Thus, \texttt{chmod 644} grants read/write permissions to the owner, and read-only permissions to group members and all other users.\\[2mm]
\textbf{Q.26 (Ans: A)} In Linux, newly created regular files start with base permissions \texttt{666} (\texttt{rw-rw-rw-}), while directories start with \texttt{777} (\texttt{rwxrwxrwx}). The \texttt{umask} bits are masked out (subtracted): for files, $666 - 022 = 644$ (\texttt{rw-r--r--}); for directories, $777 - 022 = 755$ (\texttt{rwxr-xr-x}).\\[2mm]
\textbf{Q.27 (Ans: A)} The SetUID bit on an executable file instructs the kernel to set the effective user ID (EUID) of the running process to that of the file's owner (typically root) rather than the invoking user. This allows utilities like \texttt{passwd} to safely write to the privileged \texttt{/etc/shadow} file on behalf of regular users.
\end{solutionbox}

\newpage
